% Source: https://artofproblemsolving.com/community/c3h2594794_graphs_of_functions
% Stable identity: aops-topic-2594794-q01
Let $f$ be a function.
\begin{enumerate}
\item[(a)] Show that the graph of $y=f(4-x)$ is the reflection of the graph of $y=f(x)$ in the vertical line $x=2$.
\item[(b)] A student says: ``to get $y=f(4-x)$, reflect $y=f(x)$ in the $y$-axis and then shift it $4$ units to the left.'' Is this right? Give the two correct ways of combining a reflection in the $y$-axis with a horizontal shift, and explain why the order matters.
\end{enumerate}
